\scanpage{36}
\mspage{14}
$\pi_g$ is the smallest cardinal s.th. source and target of $g$ are $\pi$-accessible.

\uline{Example} If the objects occurring as sources and targets of $g$'s in $\Phi_0$ are ``of finite presentation'', i.e. $\operatorname{Hom}(A,-)$ commutes with any filtering direct limits, we get $\pi=0$. In this case we stop the induction with $\alpha_0=$ first infinite ordinal \deleted{the first \unsure{one}}. In case the objects $A,B$ are not of finite presentation, (\unsure{still|but} they are accessible
--- we take the smallest cardinal such that $\operatorname{Hom}(A,?)$ commutes with direct limits indexed by \deleted{\unsure{ordi}} directed sets \deleted{great} \inserted{large} w.r. $\pi$.
Let \inserted{then} $\alpha_0$ be the first ordinal s.th. the set
\[
I_{\alpha_0}=\{\text{all ordinals }\alpha\text{ with }\alpha<\alpha_0\}
\]
be \deleted{great} large wrt $\pi$. We take the factorization
\diagram{\begin{tikzcd}
X\arrow[r,"i_{\alpha_0}(f)=i(f)"]\arrow[d]&X_{\alpha_0}(f)=\deleted{\Sigma}\unsure{\overline X}(f)\arrow[dl,"\Sigma_{\alpha_0}(f)=p(f)"]\\
Y&
\end{tikzcd}}
and we get $i(f)\in\Phi$, $p(f)\in\Psi$, provided we make the following assumptions;

$d)$ The \deleted{\unsure{elem}} source and target objects of the arrows of $\Phi_0$ are \uline{accessible}
(ok if $M$ a topos etc ---) and, if $\pi_0$ is the cardinality of the smallest ordinal number \inserted{$\alpha_0$} s.th. for any such $A$, $\operatorname{Hom}(A,-)$ commutes to the direct limits of type $I_{\alpha_0}$, \deleted{\unsure{we}} and \deleted{$\pi=\sup(\unsure{\pi}$}
$\pi=\sup(\pi_0,\operatorname{card}\Phi_0)$ \inserted{rather, $\sigma$}, we assume that for any filtering direct system with an indexing set $I$ of cardinality $\leq\pi$, and transition maps in $\Phi$, the direct limit exists and \deleted{$\unsure{X_n}$} $Z_i\to Z$ are equally in $\Phi$.

Then we get, not only factorization $f=pi$, with $i\in\Phi$, $p\in\Psi$, but even factorization $\uline{\overline X}(f)$ \uline{functorial} with respect to variable $f$.
\mspage{35}

